\section{The Vow of Poverty}

\textbf{Valentin Tomberg} describes the purpose of the practice of poverty:

\begin{quotex}
The practice of poverty pins down the tendencies of the thief in the human being whose male side tends to seize and female side to keep indefinitely instead of waiting for the free gift or merited fruits of one's labour.
\end{quotex}

Obviously, he does not mean material poverty in itself, since a material cause cannot effect a spiritual change. It is interesting to note how this corresponds to the Bhagavad Gita, where we read Krishna's advice to Arjuna:

\begin{quotex}
Be free from the pairs of opposites. Be always established in \emph{sattva}. Do not try to acquire what you lack or preserve what you have. Be established in the Self. (2:15)
\end{quotex}


Here he is referring to the three gunas, or fundamental forces.

\textbf{Rajas} is the active (male) principle that seeks to acquire, while \textbf{tamas} is the inertial (female) principle that seeks to horde. \textbf{Sattva}, serenity and harmony, is the equilibrium between rajas and tamas. Tomberg describes the practice of spiritual poverty in these words:

\begin{quotex}
The vow of poverty is the practice of inner emptiness that is established as a consequence of the silence of personal desires, emotions and imagination so that the soul may be capable of receiving the revelation of the Word, the life and the light from above. Poverty is the active perpetual vigil and waiting in the face of eternal sources of creativity; it is the soul ready for what is new and unexpected, it is the aptitude to learn always and everywhere, it is the \emph{condicio sine qua non} of all illumination, all revelation and all initiation.
\end{quotex}


In the light of this, our understanding of the inclinations to seize and horde is deepened. Poverty of spirit must first follow obedience. Obedience is the silencing of personal desires, emotions, and imagination, by making them subject to the dictates of reason, conscience, and legitimate authority. We may also want to add here the silencing of personal opinion.

Without this silencing, illumination and revelation are impossible; we cannot hear the new and unexpected, since all we hear are our old and persistent expectations. The desire to seize leads us to claim an understanding that is not ours. When our mind is emptied, we may receive the free gift of the Spirit which is the revelation of the Word. Otherwise, we are hearing what we want to hear, that which satisfies a secret desire, makes us ``feel'' good or important, enflames our imagination with fantasies of power or success. There is a simple touchstone we can use to discern the spirits.

\begin{enumerate}


\item Is what we hear reasonable?

\item Is what we hear good and consistent with the moral law, or does it perturb our conscience?

\item Is it consistent with spiritual and other legitimate authority?
\end{enumerate}


The second obstacle is the desire to horde. The female desire to horde means that we hang on to our personal opinions and theories. This takes many forms and it is quite common in milieus of this type. In this case, we cannot hear anything new, since everything is filtered through our prior expectations. We may belong to some political, religious, or intellectual movement that colours all our perceptions. This is subjectivity, or ``arbitrary, personal activity'', as Tomberg puts it. This subjectivity is an obstacle to the realization of the Holy Spirit. Its opposite is objectivity, as explained below:

\begin{quotex}
The prerogative of the human state is objectivity; the essential content of which is the Absolute. There is no knowledge without objectivity of the intelligence; there is no freedom without objectivity of the will; and there is no nobility without objectivity of the soul ... Esoterism seeks to realize pure and direct objectivity; this is its raison d'être.
\flright{\textsc{Frithjof Schuon}, \emph{Esoterism As Principle and As Way}}
\end{quotex}

Originally published on 23 November 2011 on medtarot.


\flrightit{Posted on 2023-03-09 by Cologero}

